\scanpage{8}
\mspage{6}

et $R=\operatorname{Spec}(\underline Z)$ (ou $\operatorname{Spec}(\underline Z_p)$, spectre des entiers p-adiques, cela reviendrait au même). Si $\underline M$ est un cristal de modules sur S, \inserted{de présentation finie, \deleted{sur}} alors en vertu de 1.5. , pour tout $s\in S$, si $s'\deleted{\unsure{a}}$ est le spectre d'une clôture parfaite de k(s), l'image inverse $\underline M(s')$ de $\underline M$ en s' peut être interprétée comme un module de type fini sur l'anneau des vecteurs de Witt W(k(s')). Ainsi, la notion de cristal de modules sur S (au sens absolu) semble convenable pour formaliser la notion de ``\uline{famille} \inserted{\uline{algébrique}''} \uline{de modules sur les anneaux de vecteurs de Witt aux différents points parfaits au dessus de} S. Bien entendu, la notion de cristal est plus fine que ça encore, et doit donner, j'espère, la bonne notion de``famille'' lorsque, disons, S est le spectre d'un corps de fonctions (donc pas nécessairement parfait). Je vais maintenant introduire une structure supplémentaire, qui \deleted{\unsure{permett|permet}} devrait permettre de formuler, de même, la notion de ``famille algébrique de modules de Dieudonné'', paramétrée par S.

Considérons le morphisme ``puissance p.ème''
\[
S\xrightarrow{\operatorname{frob}_S}S,
\]
il permet d'associer, à tout cristal (absolu) sur S, d'espèce quelconque, un cristal image inverse:
\[
\underline M^{(p)}=\operatorname{frob}_S(\underline M).
\]
On appelle p-\uline{cristal} sur S un cristal $\underline M$ sur S, muni d'un morphisme de cristaux
\[
F:\underline M^{(p)}\longrightarrow\underline M.
\]
Evidemment, les p-cristaux \inserted{sur S} d'espèce $\underline F$ donnée forment encore une catégorie, dépendant fonctoriellement de $\underline F$ (pour les foncteurs \uline{covariants} cartésiens de catégories fibrées), ce qui permet par exemple, pour les p-cristaux de modules sur S, d'introduire toutes les opérations tensorielles habituelles covariantes (produits tensoriels, puissances alternées et symétriques etc).

Pour \deleted{formuler} définir le \uline{dual} d'un p-cristal de modules, il y a lieu d'\unsure{introduite|introduire} une notions duale de celle de p-cristal d'espèce $\underline F$, c'est celle de $p^{-1}$-\uline{cristal} d'espèce $\underline F$ : c'est un cristal d'espèce $\underline F$, avec un morphisme de cristaux en sens inverse:
\margin{7}
